\section{Proposed Adaptive Domain Decomposition Method}
\label{sec:method}

\subsection{Overview of the Method}

We propose an adaptive domain decomposition framework in which the computational
domain is partitioned into subregions, each associated with a local neural network
expert. The decomposition mechanism is based on a decision tree structure, guided by
the theoretical connection between tree-based representations and smoothness
introduced in Section~\ref{sec:background}. In contrast, the global composition
of local experts follows a partition-of-unity formulation inspired by existing
domain decomposition and mixture-of-experts methods~\cite{tanno2020adaptive,hu2022apinn,moseley2021fbpinn}.

\vspace{0.5em}

\subsection{Tree Construction and Domain Decomposition}
\label{sec:tree_construction}

The goal of the decomposition is to partition the domain into regions that
reflect the local complexity of the solution. To better understand the desired
structure, we construct reference trees using ground truth solutions,
where the decomposition is determined directly from the geometric wavelet
representation. Examples for several benchmark PDEs are shown in
Figure~\ref{fig:perfect_trees}. These trees provide intuition regarding the
regions that are expected to receive additional refinement under an ideal
decomposition and serve as a reference for the adaptive procedures considered
in this work.

\begin{figure}[H]
\centering

\begin{subfigure}{0.49\textwidth}
    \centering
    \includegraphics[width=\textwidth]
    {images/section3_method/schrodinger_perfect_tree.png}
    \caption{Schrödinger}
\end{subfigure}
\hfill
\begin{subfigure}{0.49\textwidth}
    \centering
    \includegraphics[width=\textwidth]
    {images/section3_method/allen_cahn_perfect_tree.png}
    \caption{Allen--Cahn}
\end{subfigure}

\vspace{0.5em}

\begin{subfigure}{0.49\textwidth}
    \centering
    \includegraphics[width=\textwidth]
    {images/section3_method/burgers1d_perfect_tree.png}
    \caption{Burgers}
\end{subfigure}
\hfill
\begin{subfigure}{0.49\textwidth}
    \centering
    \includegraphics[width=\textwidth]
    {images/section3_method/kdv_perfect_tree.png}
    \caption{KdV}
\end{subfigure}

\caption{
Examples of reference trees constructed from ground truth solutions.
For each benchmark, the left panel shows the initial tree before pruning,
the middle panel shows the retained decomposition after pruning according
to the geometric wavelet norm, and the right panel shows the corresponding
tree hierarchy. The resulting partitions reveal regions of elevated local
complexity and illustrate the structures that the proposed adaptive
decomposition aims to recover during training.
}
\label{fig:perfect_trees}
\end{figure}

\textbf{[TODO: Consider smaller M.]}

\vspace{0.5em}

\noindent In practice, we consider two approaches for constructing the tree:

\begin{itemize}
    \item \textbf{Incremental splitting:} training is performed for a fixed number
    of epochs (or until a convergence criterion), after which a single split is
    introduced by fitting a shallow decision tree (depth 1). The split is accepted
    if it satisfies a norm-based criterion motivated by the theory. This approach may be advantageous when an initial vanilla PINN is unable to
    capture sufficient structure of the solution to reliably estimate a complete
    tree.

    \item \textbf{$M$-term tree construction:} following mixture-of-experts
training practices that advocate training components
separately~\cite{tanno2020adaptive,hu2022apinn}, the root/base model is first
trained to convergence using a standard PINN formulation. This root model
serves a dual purpose: it provides the approximation from which the geometric
wavelet tree is constructed, and it is subsequently reused as the base expert
in the final adaptive composition. A geometric wavelet tree is then
constructed from this converged approximation and compressed using an $M$-term
approximation, retaining the $M$ largest geometric wavelet coefficients.
The tree closure strategy depends on the variant:
\begin{itemize}
    \item \textbf{AToE (Additive Tree of Experts):} only ancestors of the selected
    nodes are retained; siblings are dropped. Unselected regions are carried by
    the additive root contribution (see Section~\ref{sec:global_composition}).
    \item \textbf{AToE-Leaves and ANT:} all ancestors of the selected nodes,
    together with their siblings, are retained to form a complete binary tree.
    This ensures a clean structure: in the leaves-only formulation, every spatial
    point is covered by exactly one leaf expert; in ANT, it avoids having a single
    hidden layer act as both a routing layer and an output layer.
\end{itemize}
When time-window training is employed, the global budget $M$ is distributed across the windows according to a prescribed allocation strategy (e.g., uniform, linear, or quadratic). An independent $M_i$-term tree is then constructed within each window, with $\sum_i M_i = M$.
\end{itemize}

\noindent The $M$-term construction avoids the need to manually select a wavelet-norm
threshold and is generally more robust to such choices. Furthermore, training
the root to convergence before tree construction provides a reliable
approximation for identifying the most significant structures in the geometric
wavelet representation. The converged root then serves as the base expert in
the final composition, ensuring that the global solution structure is already
well-captured before local refinements are introduced. Examples of the
resulting decompositions for several benchmark PDEs are shown in
Figure~\ref{fig:mterm_trees}.

\vspace{0.5em}

\begin{figure}[H]
\centering

\begin{subfigure}{0.4\textwidth}
\centering
\includegraphics[width=\textwidth]
{images/section3_method/mterm_schrodinger.png}
\caption{Schrödinger}
\end{subfigure}
\hfill
\begin{subfigure}{0.4\textwidth}
\centering
\includegraphics[width=\textwidth]
{images/section3_method/mterm_allen_cahn.png}
\caption{Allen--Cahn}
\end{subfigure}

\vspace{0.5em}

\begin{subfigure}{0.4\textwidth}
\centering
\includegraphics[width=\textwidth]
{images/section3_method/mterm_burgers.png}
\caption{Burgers}
\end{subfigure}
\hfill
\begin{subfigure}{0.4\textwidth}
\centering
\includegraphics[width=\textwidth]
{images/section3_method/mterm_kdv.png}
\caption{KdV}
\end{subfigure}

\textbf{[TODO: Update KS to be by window, consider smaller M.]}

\vspace{0.5em}

\caption{
Examples of adaptive decompositions obtained using the proposed $M$-term tree
construction procedure. For each benchmark, an initial vanilla PINN was first
trained for 2000 epochs using a ResNet architecture. A geometric wavelet tree
of maximum depth 15 was then constructed from the resulting approximation and
compressed using the $M$-term-by-norm procedure. The resulting trees capture many of the dominant structures
observed in the corresponding reference decompositions.
}
\label{fig:mterm_trees}
\end{figure}

\subsection{Local Neural Network Experts}

Each subdomain is associated with a neural network expert. We consider standard
architectures such as multilayer perceptrons and residual networks.

\noindent The capacity assigned to each expert may be adapted according to the
local complexity of the corresponding region. In particular, regions associated
with larger geometric wavelet norms may require larger models, while smoother
regions may be represented using smaller networks. This possibility is motivated
by the relationship between geometric wavelet coefficients, smoothness, and
approximation rates established in the previous section.

\vspace{0.5em}

\subsection{Global Composition of Local Experts}

\label{sec:global_composition}

Let $\Omega \subset \mathbb{R}^d$ denote the computational domain, and let
$\{\Omega_i\}_{i=0}^{k}$ be the set of regions obtained from the hierarchical
decision tree decomposition, where each region is an axis-aligned hypercube
\[
\Omega_i = [a_{i,1}, b_{i,1}] \times \cdots \times [a_{i,d}, b_{i,d}].
\]
Each region $\Omega_i$ is associated with a neural network expert $u_i(x,t)$.

\subsubsection{Compact Smoothstep Windows.}
To smoothly blend the contributions of different experts while ensuring
compatibility with the PDE's differential order, we define compactly-supported
window functions with configurable smoothness. Let $N$ denote the smoothness
order, chosen such that $N \geq$ the spatial order of the PDE (e.g., $N \geq 2$
for second-order PDEs like Burgers or Allen--Cahn, $N \geq 3$ for KdV).

We first define the smoothstep polynomial $S_N : [0,1] \to [0,1]$, which
satisfies $S_N(0) = 0$, $S_N(1) = 1$, and $S_N^{(k)}(0) = S_N^{(k)}(1) = 0$
for $k = 1, \ldots, N$:
\begin{align}
S_1(t) &= 3t^2 - 2t^3 \quad \text{($C^1$)} \\
S_2(t) &= 6t^5 - 15t^4 + 10t^3 \quad \text{($C^2$)} \\
S_3(t) &= 35t^4 - 84t^5 + 70t^6 - 20t^7 \quad \text{($C^3$)} \\
S_4(t) &= 126t^5 - 420t^6 + 540t^7 - 315t^8 + 70t^9 \quad \text{($C^4$)}
\end{align}

The one-sided compact ramp function is defined as:
\begin{equation}
\rho_N(s) = \begin{cases}
0 & \text{if } s \leq 0 \\
S_N(s) & \text{if } 0 < s < 1 \\
1 & \text{if } s \geq 1
\end{cases}
\end{equation}

For region $\Omega_i$ with bounds $[a_{i,j}, b_{i,j}]$ along dimension $j$,
let $\delta_{i,j} = \alpha (b_{i,j} - a_{i,j})$ denote the collar (transition)
half-width, where $\alpha > 0$ is a configurable fraction (typically $\alpha = 0.2$).
The one-dimensional window along dimension $j$ is:
\begin{equation}
\omega_{i,j}(x_j) = \rho_N\!\left(\frac{x_j - (a_{i,j} - \delta_{i,j})}{\delta_{i,j}}\right) \cdot \rho_N\!\left(\frac{(b_{i,j} + \delta_{i,j}) - x_j}{\delta_{i,j}}\right)
\label{eq:window_1d}
\end{equation}

This produces a \emph{flat-top} window: $\omega_{i,j}(x_j) = 1$ for $x_j \in [a_{i,j}, b_{i,j}]$,
smooth $C^N$ ramps in the collars $[a_{i,j} - \delta_{i,j}, a_{i,j}]$ and
$[b_{i,j}, b_{i,j} + \delta_{i,j}]$, and \emph{exactly zero} outside
$[a_{i,j} - \delta_{i,j}, b_{i,j} + \delta_{i,j}]$.

The region indicator function is the tensor product over all dimensions:
\begin{equation}
\Psi_{\Omega_i}(X) = \prod_{j=1}^{d} \omega_{i,j}(x_j)
\label{eq:soft_indicator}
\end{equation}

The expert associated with the root node (corresponding to the entire domain)
has a constant indicator $\Psi_{\Omega_0}(X) = 1$ everywhere.

\subsubsection{Normalization Strategies.}
The normalization of the indicator functions depends on the chosen global
composition strategy and varies across the three variants:

\begin{itemize}
    \item \textbf{AToE (Additive):} Per-level background normalization. For each
    level $\ell$ in the tree, the weights are normalized against a constant
    background of 1:
    \begin{equation}
    w_i(X) = \frac{\Psi_{\Omega_i}(X)}{1 + \sum_{k:\,\text{level}(k)=\ell(i)} \Psi_{\Omega_k}(X)}
    \label{eq:normalization_atoe}
    \end{equation}
    The constant 1 in the denominator serves as a ``background slot'': it keeps
    a lone child regional ($w = \Psi/(1+\Psi) < 1$) and avoids division by zero.
    
    \item \textbf{AToE-Leaves:} Normalization over all leaf nodes:
    \begin{equation}
    \tilde{\Psi}_{\Omega_i}(X) = \frac{\Psi_{\Omega_i}(X)}{\sum_{j \in \mathcal{L}} \Psi_{\Omega_j}(X)}
    \label{eq:normalization_leaves}
    \end{equation}
    where $\mathcal{L}$ denotes the complete set of leaf nodes.
    
    \item \textbf{ANT (Adaptive Neural Tree):} Leaf-frontier normalization. At
    each training stage $\ell$, normalization is computed over the current leaf
    frontier $\mathcal{F}_\ell$ (nodes that are leaves at that stage):
    \begin{equation}
    \tilde{\Psi}_{\Omega_i}(X) = \frac{\Psi_{\Omega_i}(X)}{\sum_{j \in \mathcal{F}_\ell} \Psi_{\Omega_j}(X)}
    \label{eq:normalization_ant}
    \end{equation}
    The frontier is re-evaluated as levels are added: parent nodes retire from
    the blend when their children are spawned.
\end{itemize}

\noindent In all cases, the normalized indicators satisfy the partition of unity
property $\sum_i \tilde{\Psi}_{\Omega_i}(X) = 1$ (or sum to at most 1 for AToE's
additive formulation per level).

\subsubsection{Global Solution Formulations.}
Using the normalized indicators defined above, the global solution $u_\theta(x,t)$
is constructed from the local experts in three different variants:

\begin{itemize}
    \item \textbf{Adaptive Tree of Experts (AToE).}
    In this formulation, the root/base expert contributes with coefficient
    exactly 1, and all other experts add weighted residual corrections:
    \begin{equation}
    u_\theta(x,t) = u_0(x,t) + \sum_{\ell=1}^{L} \sum_{i:\,\text{level}(i)=\ell}
    w_i(x,t)\, u_i(x,t)
    \label{eq:global_atoe}
    \end{equation}
    where $w_i$ is defined in Equation~\eqref{eq:normalization_atoe}. This
    additive structure enables coarse-to-fine refinement: the base captures
    the global solution structure, while experts at finer levels add localized
    corrections. Experts are initialized to contribute $u_i \equiv 0$ at spawn
    time, ensuring continuity of the composition.
    
    \item \textbf{AToE-Leaves.}
    Only the experts associated with leaf regions participate in the global
    solution:
    \begin{equation}
    u_\theta(x,t) =
    \sum_{j \in \mathcal{L}}
    \tilde{\Psi}_{\Omega_j}(x,t)\, u_j(x,t)
    \label{eq:global_atoe_leaves}
    \end{equation}
    The base/root is excluded from the composition once leaf experts are spawned.
    
    \item \textbf{Adaptive Neural Tree (ANT).}
    In the ANT formulation, information is propagated through the hierarchy by
    passing the activation of the last hidden layer from each parent expert to
    its children. Let $A_i(x,t)$ denote the activation of the last hidden layer
    of expert $u_i$, and let $\mathrm{parent}(i)$ denote the parent of node $i$.
    The output of expert $u_i$ is defined by
    \begin{equation}
    u_i\!\left(A_{\mathrm{parent}(i)}(x,t)\right)
    =
    \left(
    u_i^{*}(x,t),
    A_i(x,t)
    \right),
    \label{eq:ant_activation}
    \end{equation}
    where $u_i^{*}(x,t)$ denotes the prediction of the expert in the output space of the problem. The global solution
    is then given by
    \begin{equation}
    u_\theta(x,t) =
    \sum_{j \in \mathcal{F}_\ell}
    \tilde{\Psi}_{\Omega_j}(x,t)\, u_j^{*}(x,t)
    \label{eq:global_ant}
    \end{equation}
    where $\mathcal{F}_\ell$ is the current leaf frontier at training stage $\ell$.
\end{itemize}

\noindent All three variants produce a globally defined solution from the
same adaptive decomposition and support the allocation of computational
resources according to local complexity. The differences between them lie in
the interaction between experts and the resulting trade-offs between model
capacity, parameter efficiency, information sharing, training stability, and
computational complexity.

\begin{figure}[H]
\centering
\includegraphics[
    width=0.9\textwidth,
    trim={0cm 8cm 0cm 1cm},
    clip]{images/section3_method/atoe_variants}
\caption{Visualization of the AToE, AToE-Leaves, and ANT formulations.}
\label{fig:atoe_variants}
\end{figure}

\subsection{Staged Training Flow}
\label{sec:training_flow}

The training procedure follows a staged coarse-to-fine approach, where experts
are trained level-by-level with prior levels frozen. This strategy is inspired
by progressive training techniques used in mixture-of-experts
systems~\cite{tanno2020adaptive} and domain decomposition methods for
PINNs~\cite{jagtap2020xpinn,moseley2021fbpinn,hu2022apinn}.

The specific training flow depends on the variant:

\subsubsection{AToE and ANT (Staged Coarse-to-Fine).}

\begin{enumerate}
    \item \textbf{Root Phase:} Train the root/base model on the full PDE loss
    (residual + boundary conditions + initial conditions) until convergence or
    a fixed epoch budget.
    
    \item \textbf{Tree Construction:} Compute the $M$-term tree once from the
    root's prediction. Group accepted nodes by depth into levels
    $\ell = 1, \ldots, L$.
    
    \item \textbf{Per-Level Training:} For each level $\ell = 1, \ldots, L$
    (coarse to fine):
    \begin{itemize}
        \item Spawn level-$\ell$ neural network experts, initialized from their
        trained parents.
        \item Freeze the base and all lower levels ($\ell' < \ell$).
        \item Build a fresh optimizer over level-$\ell$ parameters only, with
        learning rate $\text{lr}_\ell = \text{lr}_{\ell-1} \times \gamma$, where
        $\gamma$ is a decay factor (inspired by per-subdomain learning rate
        schedules in AB-PINN).
        \item Train on the full PDE loss until per-level stopping criteria are
        met (patience-based early stopping with minimum/maximum epoch bounds).
        \item Hard-freeze level $\ell$ before proceeding.
    \end{itemize}
    
    \item \textbf{Final Joint Fine-Tuning:} Unfreeze all parameters and train
    the complete composition jointly with a lower learning rate and warmup
    period to absorb the optimizer reset.
\end{enumerate}

\subsubsection{AToE-Leaves (Joint Training).}

\begin{enumerate}
    \item \textbf{Root Phase:} Same as above.
    
    \item \textbf{Tree Construction + Spawn:} Compute the $M$-term tree and
    spawn all leaf experts in a single step.
    
    \item \textbf{Leaf Training:} Train all leaf experts jointly (base retired
    from composition) until convergence.
\end{enumerate}

\noindent The staged approach for AToE/ANT ensures that each level learns
meaningful corrections before finer levels are introduced, preventing
interference between levels during optimization. The joint approach for
AToE-Leaves is simpler but may be less stable for deep trees.

\subsection{Training and Implementation Details}

The training procedure follows standard PINN optimization, with several enhancements
to improve convergence and stability. In particular, we incorporate:

\begin{itemize}
    \item Adaptive sampling strategies~\cite{wu2023rad}.
    \item Causal training for time-dependent problems~\cite{wang2022causal}.
    \item Adaptive loss weighting (e.g., LRA)~\cite{wang2021lra}.
    \item Fourier feature embeddings~\cite{tancik2020fourier}.
    \item Weight factorization techniques~\cite{wang2021weightfactorization}.
    \item Second-stage optimizers (e.g., L-BFGS, SSBroyden)~\cite{kiyani2025optimizer}.
    \item Double precision arithmetic.
    \item Expert initialization techniques: parent-weight copying for hidden
    layers with variant-specific output layer handling (zeroed for AToE additive
    residuals, copied for AToE-Leaves/ANT partition-of-unity continuity).
\end{itemize}

\textbf{[TODO: Clarify which techniques are used in final experiments]}
